\section{C6 — 작은 실험으로 큰 실험을 맞출 수 있나}

\claimx{학습 행 수를 사다리로 세워 작은 넷에서만 곡선을 적합하고 큰 자리를 예측하면, 예측은 그 자리의 이중 붓스트랩 $SE$ 보다 크게 빗나간다}{큰 자리의 실측이 적합에 한 비트라도 들어가거나, 절대 오차가 그 자리의 $SE$ 보다 작으면 이 주장은 떨어진다.}

사다리는 생산 설정(능형 $\lambda$ 를 속겹 교차검증으로 고른다)에서 학습 행 수만 바꾼 것이다. 적합에 쓴 자리는 학습 행 1600 이하 넷이고, 나머지 자리의 실측은 적합에 들어가지 않았다. 가장 큰 자리(학습 행 $37,531$)의 실측 스피어만은 $0.328594$ 인데, 로그선형 곡선은 $0.244005$, 포화 멱법칙 곡선은 $0.198773$ 만큼 빗나간다. 그 자리의 이중 붓스트랩 $SE$ 는 $0.053097$ 다. RMSE 는 각각 $0.202308$ 와 $0.175659$ 이다.

\section{C4 — 요인을 한 모수 족으로 쓰고 예산을 맞추면 순위가 바뀐다}

\claimx{요인의 크기를 「스팬」으로 재면 그 값은 격자의 수준 수에 단조라 요인 크기의 자가 아니다. 각 요인을 한 모수 족으로 쓰고 학습 행 수를 맞춘 뒤 정규화 기울기와 같은 예산에서의 최대 $|\Delta|$ 로 재야 한다}{어떤 수준에서 학습 행 수가 예산과 다르거나, 만들어진 학습 자료의 지문이 그 요인이 움직이기로 한 성분 밖에서 움직이면 이 주장은 떨어진다.}

다섯 요인을 각각 모수 하나로 다시 썼다: 능형 $\lambda$ 의 상용로그 · 남긴 특징 개수 · 남긴 행 비율 · 도메인 가중 지수 · 원천 섞음 비율. 모든 수준에서 학습 행 수를 같게 맞췄고(예산이 맞은 수준 36 / 36), 만들어진 학습 자료의 지문으로 「한 칸만 움직였나」를 확인했다(31 / 36). 밑판은 예산을 맞춘 쪽이 $0.242961$, 예산이 없는 생산 설정이 $0.328594$ 다.

\begin{table}[h]\centering
\begin{tabular}{llll}
\hline
순위 & 요인 & 같은 예산 최대 $|\Delta|$ & 붓스트랩 1위 비율 \\
\hline
1 & R5 원천 혼합 (data mixture) & 0.546775 & 0.9975 \\
2 & R4 도메인 가중 (objective) & 0.388715 & 0.0025 \\
3 & R1 능형 $\lambda$ (optimizer) & 0.204444 & 0 \\
4 & R3 학습 창 (curriculum) & 0.156167 & 0 \\
5 & R2 특징 개수 (architecture) & 0.150615 & 0 \\
\hline
\end{tabular}
\caption{같은 예산에서의 최대 $|\Delta|$ 로 매긴 요인 순위. 스팬은 내지 않는다.}
\end{table}

전역으로 창을 자르면 원천 혼합이 따라오지만(원천 비율이 움직인 수준 6 / 7), 원천 안에서 자르면 따라오지 않는다(5 / 7). 표준오차는 학습과 유보를 함께 되뽑는 이중 붓스트랩으로 냈고, 채택 문턱 둘을 다 넘은 수준은 10 / 32 이다.
